\newpage
\section{Exhaustive Practice Problem Bank (Chapter 7)}
\label{sec:ch07_problem_bank}

\begin{tcolorbox}[enhanced,colback=subtlegreen,colframe=cengagegreen,arc=2mm,boxrule=1pt,
    title=\textbf{\large Chapter 7 Problem Bank}]
All problems sourced from: \textbf{[GRB]} pp.~37--45, \textbf{[CNG-TH]} pp.~91--110, \textbf{[NA]} pp.~21--43, \textbf{[NRJ]} Ex-I \& II, \textbf{[PRSN]} Sec.~C, \textbf{[CNG-DPP]} DPP-4--5. Full solutions in Complete Solutions Edition.
\end{tcolorbox}

% ============================================================
\subsection{7.1 Primary Batteries}
% ============================================================

\begin{problembox}
\textbf{Problem 7.1} \hfill \textbf{[GRB Ex-I Q22 / CNG-TH Concept App]}

In the dry cell (Leclanch\'{e} cell), the cathode is:
\begin{enumerate}[label=(\alph*)]
    \item Zinc container
    \item Graphite rod
    \item Ammonium chloride paste
    \item Manganese(IV) oxide alone
\end{enumerate}
\end{problembox}
\begin{solution}
\textbf{Correct: (b)}

In the dry cell, the \textbf{graphite (carbon) rod} is the cathode, surrounded by a paste of \ce{MnO2} + \ce{NH4Cl}. The cathode reaction: $\ce{2MnO2 + 2NH4+ + 2e- -> Mn2O3 + 2NH3 + H2O}$.\\
The zinc container is the anode. The carbon rod is the physical cathode contact.
\end{solution}

\begin{problembox}
\textbf{Problem 7.2} \hfill \textbf{[CNG-DPP DPP-4 Q2]}

In the mercury button cell, the net cell reaction is:
\begin{enumerate}[label=(\alph*)]
    \item \ce{Zn + HgO -> ZnO + Hg}
    \item \ce{Hg + ZnO -> ZnHgO}
    \item \ce{Zn + Hg -> ZnHg (amalgam)}
    \item \ce{ZnO + Hg -> Zn + HgO}
\end{enumerate}
\end{problembox}
\begin{solution}
\textbf{Correct: (a)}

Anode: $\ce{Zn + 2OH- -> ZnO + H2O + 2e-}$\\
Cathode: $\ce{HgO + H2O + 2e- -> Hg + 2OH-}$\\
Net: $\ce{Zn + HgO -> ZnO + Hg}$

Both electrode products (ZnO, Hg) are solids/liquids of constant activity, hence the voltage remains very stable throughout discharge --- the key advantage of this cell.
\end{solution}

\begin{problembox}
\textbf{Problem 7.3} \hfill \textbf{[NRJ Ex-I Q18]}

\textbf{[Assertion--Reason]}

\textbf{Statement-1:} The alkaline battery has a longer shelf life than the Leclanch\'{e} (dry) cell.

\textbf{Statement-2:} In an alkaline battery, the electrolyte is \ce{KOH}(aq), which prevents corrosion of the zinc anode by reducing the rate of local cell action at impurity sites.

\begin{enumerate}[label=(\alph*)]
    \item Both are true; S-2 is the correct explanation.
    \item Both are true; S-2 is NOT the correct explanation.
    \item S-1 is true; S-2 is false.
    \item S-1 is false; S-2 is true.
\end{enumerate}
\end{problembox}
\begin{solution}
\textbf{Correct: (a)}

Both statements are true. The alkaline electrolyte (\ce{KOH}) reduces local corrosion of Zn (which occurs readily in acidic \ce{NH4Cl} environment due to local galvanic cells at impurity sites), gives more stable voltage under load, and provides greater total capacity. S-2 correctly explains the mechanism behind S-1.
\end{solution}

% ============================================================
\subsection{7.2 Lead-Acid Battery --- Quantitative}
% ============================================================

\begin{problembox}
\textbf{Problem 7.4} \hfill \textbf{[NA Level-2 Q22 / CNG-TH Illustration 9]}

During discharge of a lead-acid battery, the density of sulfuric acid decreases from $1.28\,\text{g\,mL}^{-1}$ to $1.10\,\text{g\,mL}^{-1}$. The cell delivers $10\,\text{A}$ for $t$ hours. If the total mass of \ce{H2SO4} consumed per cycle is $29.4\,\text{g}$, find $t$.

[$M_{\ce{H2SO4}} = 98$; for the cell reaction: $\ce{Pb + PbO2 + 2H2SO4 -> 2PbSO4 + 2H2O}$; $n=2$ per mole of cell reaction]
\end{problembox}
\begin{solution}
Moles of \ce{H2SO4} consumed: $29.4/98 = 0.3\,\text{mol}$.

From the net reaction: 2~mol \ce{H2SO4} are consumed per 2~mol $e^-$ transferred ($n=2$ per formula unit).\\
Moles of electrons = moles of \ce{H2SO4} (1:1 ratio above):
Actually from: $\ce{Pb + PbO2 + 2H2SO4 -> 2PbSO4 + 2H2O}$ with $n=2$:
$0.3\,\text{mol}\,\ce{H2SO4} \to 0.3/2 = 0.15\,\text{mol}$ of cell reaction $\to 0.15 \times 2 = 0.3\,\text{mol}\,e^-$\\
$Q = 0.3 \times 96500 = 28950\,\text{C}$\\
$t = Q/I = 28950/(10 \times 3600) = \mathbf{0.804\,\text{h} \approx 48.3\,\text{min}}$
\end{solution}

\begin{problembox}
\textbf{Problem 7.5} \hfill \textbf{[GRB Numerical Q8 / NRJ Ex-II Q24]}

In a lead-acid battery, during charging, what happens to the density of the sulfuric acid solution? Also, write the net charging reaction.
\end{problembox}
\begin{solution}
During \textbf{charging}, the reverse cell reaction occurs:
\[\ce{2PbSO4(s) + 2H2O(l) -> Pb(s) + PbO2(s) + 2H2SO4(aq)}\]
\ce{H2SO4} is \textbf{regenerated} $\to$ concentration increases $\to$ \textbf{density increases}.

This is why measuring the density (specific gravity) with a hydrometer indicates the state of charge:
\begin{itemize}
    \item Fully charged: $\rho \approx 1.28\,\text{g\,mL}^{-1}$
    \item Fully discharged: $\rho \approx 1.10\,\text{g\,mL}^{-1}$
\end{itemize}
\end{solution}

\begin{problembox}
\textbf{Problem 7.6} \hfill \textbf{[NA Level-3 Q15]}

\textbf{[Matrix Match]} Match the cell with its electrode reaction:

\begin{tabular}{ll@{\quad}ll}
\toprule
\textbf{Column I (Cell)} & & \textbf{Column II (Anode Reaction)} & \\
\midrule
(A) Dry cell & & (p) $\ce{Zn + 2OH- -> ZnO + H2O + 2e-}$ & \\
(B) Mercury cell & & (q) $\ce{Zn -> Zn^2+ + 2e-}$ & \\
(C) Lead-acid (discharge) & & (r) $\ce{Pb + SO4^2- -> PbSO4 + 2e-}$ & \\
(D) Nicad (discharge) & & (s) $\ce{Cd + 2OH- -> Cd(OH)2 + 2e-}$ & \\
\bottomrule
\end{tabular}
\end{problembox}
\begin{solution}
(A) $\to$ (q): Dry cell anode = Zn metal in \ce{NH4Cl} (mildly acidic): $\ce{Zn -> Zn^2+ + 2e-}$\\
(B) $\to$ (p): Mercury cell anode = Zn in alkaline KOH: $\ce{Zn + 2OH- -> ZnO + H2O + 2e-}$\\
(C) $\to$ (r): Lead-acid discharge anode: $\ce{Pb + SO4^2- -> PbSO4 + 2e-}$\\
(D) $\to$ (s): Nicad discharge anode: $\ce{Cd + 2OH- -> Cd(OH)2 + 2e-}$
\end{solution}

% ============================================================
\subsection{7.3 Fuel Cells and Efficiency}
% ============================================================

\begin{problembox}
\textbf{Problem 7.7} \hfill \textbf{[ATK Topic 6C / NRJ Ex-II Q28]}

For the alkaline hydrogen--oxygen fuel cell, the overall reaction is $\ce{H2(g) + \frac{1}{2}O2(g) -> H2O(l)}$ with $\Delta H^\circ = -286\,\text{kJ\,mol}^{-1}$ and $E^\circ_{\text{cell}} = 1.23\,\text{V}$. Calculate the thermodynamic efficiency of this fuel cell.
\end{problembox}
\begin{solution}
$n = 2$ electrons transferred (per mol of \ce{H2}).\\
$\Delta G^\circ = -nFE^\circ = -2 \times 96500 \times 1.23 = -237.4\,\text{kJ\,mol}^{-1}$\\
$\eta = \frac{|\Delta G^\circ|}{|\Delta H^\circ|} = \frac{237.4}{286} = 0.830 = \mathbf{83.0\%}$

Compare with a Carnot engine operating between $T_H = 1000\,\text{K}$ and $T_C = 300\,\text{K}$: $\eta_{\text{Carnot}} = 1 - 300/1000 = 70\%$.
The fuel cell surpasses even this ideal Carnot efficiency because it is not a heat engine.
\end{solution}

\begin{problembox}
\textbf{Problem 7.8} \hfill \textbf{[CNG-TH Advanced Q / GRB Ex-I Q29]}

In a PEM fuel cell, protons travel from:
\begin{enumerate}[label=(\alph*)]
    \item Cathode through external circuit to anode
    \item Anode through the Nafion membrane to cathode
    \item Anode through the external circuit to cathode
    \item Cathode through the Nafion membrane to anode
\end{enumerate}
\end{problembox}
\begin{solution}
\textbf{Correct: (b)}

In a PEM fuel cell:
\begin{itemize}
    \item At the anode: $\ce{H2 -> 2H+ + 2e-}$
    \item \ce{H+} (protons) travel through the \textbf{Nafion polymer membrane} from anode to cathode (the membrane is a selective proton conductor).
    \item Electrons travel through the external circuit (generating useful current).
    \item At cathode: $\ce{O2 + 4H+ + 4e- -> 2H2O}$
\end{itemize}
\end{solution}

% ============================================================
\subsection{7.4 Corrosion}
% ============================================================

\begin{problembox}
\textbf{Problem 7.9} \hfill \textbf{[GRB Ex-I Q31 / CNG-DPP DPP-5 Q3]}

Rusting of iron requires the simultaneous presence of:
\begin{enumerate}[label=(\alph*)]
    \item Oxygen only
    \item Water only
    \item Both oxygen and water
    \item Carbon dioxide only
\end{enumerate}
\end{problembox}
\begin{solution}
\textbf{Correct: (c)}

Rusting is an electrochemical process. Both water (as electrolyte/medium) and oxygen (as cathodic oxidant) are required. Pure dry oxygen or pure water alone does not cause rusting. Salt (electrolyte) accelerates rusting by increasing ionic conductivity.
\end{solution}

\begin{problembox}
\textbf{Problem 7.10} \hfill \textbf{[NA Level-2 Q30 / NRJ Ex-I Q22]}

A piece of iron is connected to a piece of copper by a wire and immersed in salt water. Which of the following is correct?

\begin{enumerate}[label=(\alph*)]
    \item Iron acts as cathode and is protected
    \item Copper corrodes because it is less active
    \item Iron acts as anode and corrodes preferentially
    \item Both metals corrode at equal rates
\end{enumerate}
\end{problembox}
\begin{solution}
\textbf{Correct: (c)}

$E^\circ(\ce{Fe^2+/Fe}) = -0.44\,\text{V}$ and $E^\circ(\ce{Cu^2+/Cu}) = +0.34\,\text{V}$.\\
Iron has more negative $E^\circ$, so it is more easily oxidized (more active). Iron acts as the \textbf{anode}: $\ce{Fe -> Fe^2+ + 2e-}$. Copper acts as the cathode: $\ce{O2 + 2H2O + 4e- -> 4OH-}$. Iron corrodes; copper is protected. This is galvanic corrosion, where the less noble metal always acts as the anode.
\end{solution}

\begin{problembox}
\textbf{Problem 7.11} \hfill \textbf{[CNG-TH Q / GRB Ex-I Q34]}

\textbf{[Multiple Correct]} Which of the following methods can prevent corrosion of iron?

\begin{enumerate}[label=(\alph*)]
    \item Coating with zinc (galvanizing)
    \item Connecting to a copper wire
    \item Connecting to a magnesium block
    \item Applying cathodic protection with impressed current
\end{enumerate}
\end{problembox}
\begin{solution}
\textbf{Correct: (a), (c), (d)}

(a) Galvanizing: Zn is more active, acts as sacrificial anode. \checkmark\\
(b) Connecting to Cu: Cu is less active; Fe acts as anode --- corrosion \textit{accelerates}. \texttimes\\
(c) Magnesium block: Mg ($E^\circ = -2.37\,\text{V}$) is more active than Fe ($-0.44\,\text{V}$); Mg sacrifices itself. \checkmark\\
(d) Impressed current cathodic protection: makes Fe the cathode; prevents anodic dissolution. \checkmark
\end{solution}

\begin{problembox}
\textbf{Problem 7.12} \hfill \textbf{[NRJ Ex-II Q32 / PRSN Ex-II Q17]}

Write the overall reaction for rusting of iron in neutral water and calculate the standard EMF of the corrosion cell. [$E^\circ(\ce{Fe^2+/Fe}) = -0.44\,\text{V}$; $E^\circ(\ce{O2/OH-}) = +0.40\,\text{V}$]
\end{problembox}
\begin{solution}
\textbf{Anode (oxidation):} $\ce{Fe -> Fe^2+ + 2e-}$ ($\times2$)\\
\textbf{Cathode (reduction):} $\ce{O2 + 2H2O + 4e- -> 4OH-}$\\
\textbf{Net:} $\ce{2Fe + O2 + 2H2O -> 2Fe^2+ + 4OH- -> 2Fe(OH)2}$\\
(followed by further oxidation to \ce{Fe(OH)3} and dehydration to \ce{Fe2O3.xH2O} --- rust)

$E^\circ_{\text{cell}} = E^\circ_{\text{cathode}} - E^\circ_{\text{anode}} = +0.40 - (-0.44) = \mathbf{+0.84\,\text{V}}$

Since $E^\circ > 0$, $\Delta G < 0$: corrosion is thermodynamically spontaneous.
\end{solution}

\begin{problembox}
\textbf{Problem 7.13} \hfill \textbf{[ATK Physical Chemistry / ESS Advanced]}

\textbf{[Assertion--Reason]}

\textbf{Statement-1:} Aluminum resists corrosion in air, despite having a very negative standard reduction potential ($E^\circ = -1.66\,\text{V}$).

\textbf{Statement-2:} Aluminum forms a thin, adherent, self-repairing oxide layer (\ce{Al2O3}) on its surface that acts as a passivating film, preventing further oxidation.

\begin{enumerate}[label=(\alph*)]
    \item Both true; S-2 is the correct explanation.
    \item Both true; S-2 is NOT the correct explanation.
    \item S-1 true; S-2 false.
    \item S-1 false; S-2 true.
\end{enumerate}
\end{problembox}
\begin{solution}
\textbf{Correct: (a)}

Both statements are correct, and S-2 correctly explains S-1 via passivation. The Pilling--Bedworth ratio for \ce{Al2O3} ($R \approx 1.28$) ensures the oxide film is protective, adherent, and self-sealing. This is the basis of anodizing, which artificially thickens this layer. The same principle applies to Cr in stainless steel and Ti in titanium alloys.
\end{solution}

\begin{problembox}
\textbf{Problem 7.14} \hfill \textbf{[NA Level-3 Q18 / CNG-TH Advanced]}

\textbf{[Comprehension]} A ship hull is made of iron. To prevent corrosion, zinc blocks are attached to the hull below the water line.

(A) What role does the zinc block play?
\begin{enumerate}[label=(\alph*)]
    \item It acts as a cathodic protection using impressed current.
    \item It acts as a sacrificial anode.
    \item It acts as a barrier coating.
    \item It increases the electrical resistance of the hull.
\end{enumerate}

(B) The zinc blocks need to be replaced periodically. Why?
\begin{enumerate}[label=(\alph*)]
    \item They become contaminated with rust from the iron.
    \item They are gradually consumed (corrode away) while protecting the iron.
    \item The seawater dissolves them chemically without electrochemical action.
    \item They passivate after prolonged exposure.
\end{enumerate}
\end{problembox}
\begin{solution}
(A) \textbf{(b) Sacrificial anode.}

$E^\circ(\ce{Zn^2+/Zn}) = -0.76\,\text{V} < E^\circ(\ce{Fe^2+/Fe}) = -0.44\,\text{V}$, so Zn is more active, acts as the anode of the local galvanic cell, and oxidizes preferentially, protecting Fe (the cathode).

(B) \textbf{(b) They are gradually consumed.}

As Zn acts as the sacrificial anode, it undergoes $\ce{Zn -> Zn^2+ + 2e-}$ and is progressively dissolved into the seawater. The zinc blocks must be replaced when significantly depleted; otherwise, the iron hull loses protection and begins to corrode.
\end{solution}
